\chapter{Introduction}

\section{Context}

The healthcare sector is undergoing a significant transformation driven the Internet of Things (IoT) technologies.
Continuous, real-time monitoring of patient vital signs and ambient room conditions is a foundational requirement in any clinical environment, enabling medical staff to detect deterioration early, respond rapidly, and reduce the administrative overhead associated with manual observation rounds.

This project proposes a complete, end-to-end hospital monitoring solution composed
of four integrated subsystems:

\begin{itemize}
	\item \textbf{Embedded acquisition node (ESP32):} real-time sensor acquisition and authenticated cloud transmission.
	\item \textbf{Backend server:} database management, alert rule evaluation, and
	      user authentication.
	\item \textbf{Web dashboard:} real-time data visualisation, historical charts,
	      and alert management for medical staff.
	\item \textbf{AI prediction layer:} anomaly detection on patient vitals and
	      environmental condition assessment using cloud-hosted machine learning models.
\end{itemize}

\section{Problem Statement}

Hospital rooms require continuous monitoring of two categories of parameters:
environmental conditions (temperature and humidity), which directly affect patient
comfort, infection risk, and equipment operation; and patient vital signs (heart rate
and blood oxygen saturation), which are primary indicators of physiological status.

The current approach in many low-resource clinical settings relies on periodic manual
measurement, which introduces observation gaps and increases nursing workload.

\section{Objectives}

The project objectives are divided into functional and technical categories.

\subsection*{Functional Objectives}

\begin{enumerate}
	\item Continuously acquire temperature (°C), humidity (\%), heart rate (BPM), and
	      blood oxygen saturation (SpO\textsubscript{2}, \%) from dedicated sensors.
	\item Transmit sensor data to a cloud database at a near-real-time rate of up to 1\,Hz.
	\item Implement a web dashboard for real-time data visualisation, history browsing,
	      and alert management.
	\item Generate automated predictions and anomaly alerts via a cloud AI model.
\end{enumerate}

\subsection*{Technical Objectives}

\begin{enumerate}
	\item Availability: the system shall operate continuously with network recovery mech-
	      anisms.
	\item Reliability: invalid or inconsistent data shall be detected and flagged.
	\item Security: secure authentication and encrypted communication (HTTPS).
	\item Performance: near real-time data updates based on sampling frequency.
\end{enumerate}

\section{Contributions}

This project is the result of a full-stack collaborative effort across four subsystems.
The concrete contributions of the group are as follows:

\begin{description}
	\item[Embedded acquisition node (ESP32).]
	      Complete firmware design and implementation: dual-core FreeRTOS task architecture,
	      DHT22 and MAX30102 sensor drivers, PPG-based heart rate and SpO\textsubscript{2}
	      computation, delta-threshold upload filter, LittleFS-backed configuration portal,
	      NVS credential persistence, and automatic Wi-Fi recovery. Hardware assembly
	      including resolution of the MAX30102 pull-up resistor incompatibility.

	\item[Backend server.]
	      Design and deployment of a Node.js/tRPC service layer with Firebase Admin SDK
	      integration, Role-Based Access Control (RBAC), Zod schema validation, a hybrid
	      alerting engine combining fixed clinical thresholds with Z-score anomaly
	      detection, and automated SMTP email notifications.

	\item[Web dashboard.]
	      Implementation of a React/Vite single-page application with real-time Firebase
	      RTDB subscriptions, Recharts-based vital sign visualisation, an AI Insights
	      panel, and a WebWorker-based face-api.js pipeline for authorized personnel
	      tracking using an external USB camera.

	\item[AI prediction layer.]
	      Server-side implementation of a statistical anomaly detection engine (Z-score
	      over a 20-point sliding window) and a 5-minute vital sign forecasting model
	      (Ordinary Least Squares linear regression), exposed as tRPC procedures and
	      integrated into the dashboard's clinical status display.
\end{description}

\section{Organization}

The remainder of this report is structured as follows.
Chapter~2 provides a detailed problem analysis, defining use cases, system inputs and
outputs, and platform constraints.
Chapter~3 presents the global and functional system architecture, including the hardware
and software decomposition.
Chapter~4 describes the design methodology, covering the mathematical models,
pseudocode for the main algorithms, and the development workflow.
Chapter~5 details the implementation of all subsystems, including hardware assembly
challenges, software modules, and validation procedures.
Chapter~6 reports the experimental results and quantitative validation for each
subsystem.
Chapter~7 discusses the validity of the results, system limitations, and identified
improvements.
Chapter~8 concludes the report with a synthesis of achievements and directions for
future work.
Chapter~9 provides an expanded treatment of the AI prediction layer, including its
quantitative evaluation and performance constraints.
